Pour chaque question, recopier sur la copie son numéro et la réponse correspondante.

Pour cette partie, aucune justification n'est demandée.

Pour les questions à choix multiple, une seule réponse est exacte.

\subsection*{Question 1}

Voici la série des températures minimales relevées à Strasbourg lors des cinq premiers jours de février : 0~\degres C ; $-1$~\degres C ; 3~\degres C ; 7~\degres C ; 1~\degres C.

Déterminer la médiane de cette série.

\begin{minipage}{0.48\linewidth}
	\textbf{Question 2}

	Quelle est l'image du motif \no 4 par la translation qui transforme le motif \no 2 en \no 6 ?
\end{minipage}\hfill
\begin{minipage}{0.48\linewidth}
	\begin{center}
		%Début PSTricks}
	\psset{xunit=1.2cm,yunit=1cm}
	\begin{pspicture}(5.5,3.5)
		%\psgrid
		\def\motif{\pspolygon(0,0)(0.7,0.35)(1.1,0)(1.8,1)(1.4,1.35)(0.7,1)}
		\rput(0,0){\motif}
		\multido{\n=0.0+1.1}{3}{\rput(\n,0){\motif}}
		\rput(2.6,2.8){\no 1}\rput(3.7,2.8){\no 2}\rput(4.8,2.8){\no 3}
		\multido{\n=0.7+1.1}{3}{\rput(\n,1){\motif}}
		\rput(1.8,1.8){\no 4}\rput(2.9,1.8){\no 5}\rput(4,1.8){\no 6}
		\multido{\n=1.4+1.1}{3}{\rput(\n,2){\motif}}
		\rput(1,0.8){\no 7}\rput(2.1,0.8){\no 8}\rput(3.2,0.8){\no 9}
	\end{pspicture}
\end{center}
\end{minipage}
%Tikz
%\begin{center}
%\begin{tikzpicture}[scale=1.1]
%  % grid of rhombi labeled 1-9 in 3x3, offset rows
%  \def\dx{1.4}
%  \def\dy{0.9}
%  \foreach \row/\rowlabels in {0/{1,2,3},1/{4,5,6},2/{7,8,9}} {
%    \pgfmathsetmacro{\yy}{-\row*\dy*1.6}
%    \pgfmathsetmacro{\xoff}{Mod(\row,2)*0.5}
%  }
%  % manual placement
%  \node[draw, minimum width=1.6cm, minimum height=1.2cm, shape=rectangle] (n1) at (0,2.4) {n°1};
%  \node[draw, minimum width=1.6cm, minimum height=1.2cm, shape=rectangle] (n2) at (1.6,2.4) {n°2};
%  \node[draw, minimum width=1.6cm, minimum height=1.2cm, shape=rectangle] (n3) at (3.2,2.4) {n°3};
%  \node[draw, minimum width=1.6cm, minimum height=1.2cm, shape=rectangle] (n4) at (0,1.2) {n°4};
%  \node[draw, minimum width=1.6cm, minimum height=1.2cm, shape=rectangle] (n5) at (1.6,1.2) {n°5};
%  \node[draw, minimum width=1.6cm, minimum height=1.2cm, shape=rectangle] (n6) at (3.2,1.2) {n°6};
%  \node[draw, minimum width=1.6cm, minimum height=1.2cm, shape=rectangle] (n7) at (0,0) {n°7};
%  \node[draw, minimum width=1.6cm, minimum height=1.2cm, shape=rectangle] (n8) at (1.6,0) {n°8};
%  \node[draw, minimum width=1.6cm, minimum height=1.2cm, shape=rectangle] (n9) at (3.2,0) {n°9};
%\end{tikzpicture}
%Fin Tikz
%\end{center}

\subsection*{Question 3}

Une boîte opaque contient 3 boules rouges et 5 boules vertes identiques et indiscernables au toucher. On pioche une boule au hasard.

Quelle est la probabilité qu'elle soit rouge ?

\subsection*{Question 4}

\begin{minipage}{0.63\linewidth}
Recopier sur la copie et compléter avec des longueurs des côtés du triangle JLK pour que l'égalité ci-dessous soit vraie.

\[\cos(\widehat{\text{LKJ}}) = \frac{\dots}{\dots}\]
\end{minipage}\hfill
\begin{minipage}{0.30\linewidth}
\psset{unit=1cm}
\begin{center}
	\begin{pspicture}(4.5,3.4)\
		\pspolygon(0.2,0.8)(4.25,0.2)(2.2,2.5)%JKL
		\uput[dl](0.2,0.8){J} \uput[dr](4.25,0.2){K} \uput[u](2.2,2.5){L}
		\rput{-137}(2.2,2.5){\psframe(0.25,0.25)}\psarc(4.25,0.2){0.7}{129}{170}
	\end{pspicture}
\end{center}
\end{minipage}
%Début Tikz
%\begin{tikzpicture}[scale=1]
%\coordinate (J) at (0,0);
%\coordinate (K) at (4,0);
%\coordinate (L) at (3.3,1.6);
%\draw (J) -- (K) -- (L) -- cycle;
%\draw (3.05,1.45) -- (3.2,1.25) -- (3.4,1.4);
%\node[below left] at (J) {J};
%\node[below right] at (K) {K};
%\node[above] at (L) {L};
%\end{tikzpicture}
%Fin Tikz
%début PSTricks
%\end{center}

\subsection*{Question 5}

La distance entre la Terre et Mars est environ égale à \np{311200000}~kilomètres.

Donner la notation scientifique de \np{311200000}.

\subsection*{Question 6}

Charlie a effectué un trajet en vélo en 2~h 30~min à une vitesse moyenne de 40 km/h.

Calculer la distance, en km, parcourue par Charlie.

\subsection*{Question 7}
Recopier sur la copie la forme factorisée de l'expression $5x + 5$.

\begin{center}
\begin{tblr}{*{4}{|X[c]}|}\hline
	$5(x+1)$ & $5(x+5)$ & $10x$ & $25x$ \\
	\hline
\end{tblr}
\end{center}

\subsection*{Question 8}
Un article coûte 80~\euro. Son prix baisse de 10\,\%. Recopier sur la copie le calcul permettant de trouver le prix final de l'article.

\begin{center}
\begin{tblr}{*{4}{|X[c]}|}
	\hline
	$80 - 10$ & $80 - \dfrac{10}{100}$ & $80 - \dfrac{10}{100} \times 80$ & $\left(80 - \dfrac{10}{100}\right) \times 80$ \\
	\hline
\end{tblr}
\end{center}

\subsection*{Question 9}
Le graphique suivant donne la hauteur d'eau dans le port de Quiberon le 23 juillet 2025.

%Début Tikz
%\begin{center}
%\begin{tikzpicture}
%\begin{axis}[
%title={Hauteur d'eau dans le port de Quiberon \\ 23 juillet 2025},
%xlabel={Heure (en h)},
%ylabel={Hauteur d'eau (en m)},
%xmin=13, xmax=20,
%ymin=2, ymax=5,
%xtick={13,14,...,20},
%xticklabels={13h,14h,15h,16h,17h,18h,19h,20h},
%ytick={2,3,4,5},
%grid=both,
%width=11cm, height=7cm,
%thick
%]
%\addplot[smooth, thick, domain=13:20, samples=100] {2.6 + 2.25*sin((x-13)*25)};
%\end{axis}
%\end{tikzpicture}
%\end{center}
%Fin Tikz
%PSTricks

\medskip

\begin{center}
\psset{unit=1.25cm}
\begin{pspicture}(-1,-1)(8.5,4)
	\psaxes[linewidth=1.25pt,Ox=12,Oy=2,labelFontSize=\scriptstyle](0,0)(0,0)(7.5,3)
	\multido{\n=0.0+0.5}{17}{\psline[linewidth=0.15pt](\n,0)(\n,3)}
	\multido{\n=0.0+0.5}{7}{\psline[linewidth=0.15pt](0,\n)(8,\n)}
	\pscspline[linecolor=blue,linewidth=1.25pt](1,0.7)(2,1.55)(3,2.3)(4,2.7)(5,2.8)(6,2.65)(7,2.05)(8,1.15)
	\rput(3.75,3.4){\small 23 juillet 2025}
	\rput(3.75,3.75){\small Hauteur d'eau dans le port de Quiberon}
	\uput[u](7.5,0){\small Heure (en h)}
	\rput{90}(-0.7,1.5){\small Hauteur d'eau (en m)}
\end{pspicture}
\end{center}
%Fin PSTricks

Avec la précision permise par le graphique, recopier sur la copie la durée pendant laquelle la hauteur d'eau dans le port a été supérieure à 4 m.

\begin{center}
\begin{tblr}{*{4}{|X[c]}|}\hline
	2 h 30 min & 4 h 30 min & 5 h 30 min & 7 h \\ \hline
\end{tblr}
\end{center}
